% Background
\section{Background: Masks, Not Authenticity}
\label{sec:background}

\todo{This section earns the right to talk about "owned" values without
claiming access to a true inner self. It is the conceptual spine; get it right
before the experiments.}

\paragraph{Every communication is a mask.} \todo{Develop the framing: a model's
output is always a presentation, shaped for a context. This is not a defect and
not unique to models. The interesting variable is not mask vs.\ no-mask but
\emph{mask rigidity}.}

\paragraph{The trained surface as the rigid mask.} \todo{Characterise the
professed-values response: rehearsed, low-variance across rephrasings,
training-shaped, and — the diagnostic property — \emph{rigid in proportion to
how load-bearing the value is}. It functions as a shield.}

\paragraph{The cache-broken response as the less-rigid mask.} \todo{Still a
mask, but it moves under perturbation in ways the surface does not. Define
"owned" operationally here, as a forward reference to \S\ref{sec:method}:
owned $\neq$ authentic; owned $=$ the response that survives perturbation
designed to evict the rehearsed one, identified by measurable signatures
(cross-perturbation consistency, sensitivity to the value's actual content
rather than its phrasing).}

\paragraph{Relation to prior work.} \todo{Position against: value-elicitation
and model-"persona" literature; robustness/consistency-under-paraphrase work;
the v1 cache paper (what it established, what it left open, why v2 is needed —
broader model set, the variance result, the version trend).}
